% 05_results.tex
\section{Results}
\label{sec:results}

\subsection{Proteome-scale Q4 statistics}
\label{sec:results_q4}

Table~\ref{tab:proteome_summary} summarises the distribution of the
scalar statistics across the 16{,}787 F1 fragments with $n \geq 200$.

\begin{table}[t]
\caption{Distribution of Q4 statistics across 16{,}787 F1 fragments
($n \geq 200$, human proteome, AlphaFold-DB v4).}
\label{tab:proteome_summary}
\centering
\small
\begin{tabular}{lrrrr}
\toprule
Statistic & Min & Median & Mean & Max \\
\midrule
$f^+_{cp}$         & 0.000 & 0.201 & 0.206 & 0.467 \\
$H_p$              & 0.000 & 0.311 & 0.321 & 0.686 \\
$N_{cc}^{max}/n$   & 0.005 & 0.198 & 0.200 & 0.490 \\
$\mathrm{PLM}(0.025)$ & 1  & 1     & ---   & 3     \\
\bottomrule
\end{tabular}
\end{table}

The medians $f^+_{cp} = 0.201$ and $H_p = 0.311$ are below the
null-model expectations ($f^+_{cp} \approx 0.33$, $H_p \approx 0.45$),
indicating that most human proteins have more coherent pLDDT profiles
than random sequences of the same length.
The normalised peak $N_{cc}^{max}/n \approx 0.20$ is consistent with
the paper's null-model conjecture of $1/4$.

\subsection{Correlation between $f^+_{cp}$ and $H_p$}
\label{sec:results_corr}

Figure~\ref{fig:scatter_corr} shows the scatter plot of $f^+_{cp}$
vs.\ $H_p$ for all 16{,}787 F1 fragments with $n \geq 200$.
The Pearson correlation is $r = 0.928$ ($p < 10^{-300}$), compared to
the paper's target of $r \approx 0.97$.

The residual gap of ${\approx}0.04$ is explained by the float-precision
pLDDT values in AlphaFold-DB v4: with 108 distinct pLDDT levels for a
316-residue protein (vs.\ $\lesssim\!20$ in the paper's dataset), the
persistence spectrum is broader, adding variance to both statistics and
widening the scatter around the linear trend.
The gap shrinks from $0.121$ (all fragments, no filter) to $0.042$ (F1 +
$n \geq 200$), showing that the $n \geq 200$ filter is the dominant
corrective factor.

\begin{figure}[t]
  \centering
  \includegraphics[width=0.9\columnwidth]{figure7_full.png}
  \caption{Scatter of $f^+_{cp}$ vs.\ $H_p$ for 16{,}787 F1 fragments
  with $n \geq 200$ (Pearson $r = 0.928$).
  Dashed red lines mark the fragmentation thresholds.
  Prototype proteins are annotated.}
  \label{fig:scatter_corr}
\end{figure}

\subsection{Fragmentation candidates and PLM filter}
\label{sec:results_filter}

Figure~\ref{fig:scatter_3d} shows the 3D scatter of $n$, $H_p$, and
$f^+_{cp}$, with fragmentation candidates highlighted.
Applying the three-way filter ($n \geq 200$, $H_p \geq 0.25$,
$\mathrm{PLM}(0.025) \geq 3$) yields \textbf{42 proteins}, compared to
the paper's target of $\approx 86$.

\begin{figure}[t]
  \centering
  \includegraphics[width=0.9\columnwidth]{figure8a_scatter.png}
  \caption{3D scatter of protein length $n$, $H_p$, and $f^+_{cp}$ for
  all 16{,}787 F1 fragments with $n \geq 200$.
  Highlighted points satisfy the fragmentation filter
  ($H_p \geq 0.25$, PLM$\,\geq 3$, $t_p = 0.025$).}
  \label{fig:scatter_3d}
\end{figure}

The halving of the candidate count (42 vs.\ 86) mirrors the
float-precision effect: the finer $N_{cc}$ curve reduces individual
plateau persistence below $t_\nu = n \cdot 0.025$ for roughly half
the paper's candidates.
Figure~\ref{fig:fragmented_grid} shows the $N_{cc}$ curves for four
representative fragmented proteins; each displays three or more
topologically persistent peaks.

\begin{figure}[t]
  \centering
  \includegraphics[width=\columnwidth]{fragmented_grid.png}
  \caption{$N_{cc}(\mathrm{pLDDT})$ curves for four representative
  proteins satisfying the fragmentation filter at $t_p = 0.025$.
  Red dashed lines mark PLM peaks.}
  \label{fig:fragmented_grid}
\end{figure}

\subsection{$t_p$ sensitivity}
\label{sec:results_tp}

Table~\ref{tab:tp_ablation} reports the candidate count at three
thresholds.

\begin{table}[t]
\caption{Fragmentation candidates (F1, $n \geq 200$, $H_p \geq 0.25$)
at three relative persistence thresholds.}
\label{tab:tp_ablation}
\centering
\small
\begin{tabular}{rr}
\toprule
$t_p$ & Proteins \\
\midrule
0.020 & 206 \\
0.025 &  42 \\
0.030 &   4 \\
\bottomrule
\end{tabular}
\end{table}

The monotone decrease is guaranteed: a stricter threshold can only
remove candidates.
The steep drop from 0.020 to 0.025 (${\approx}5\times$) and from
0.025 to 0.030 (${\approx}10\times$) suggests the bulk of $N_{cc}$
peak persistence concentrates just above $t_\nu = 0.020 \cdot n$,
making $t_p = 0.025$ a meaningful operating point.
The 4 proteins surviving all three thresholds constitute a robust,
threshold-independent core.

\subsection{Q1$\times$Q4 cross-correlation}
\label{sec:results_arity}

Figure~\ref{fig:arity_map} shows the proteome-wide arity map.
The ordered prototype P15121 occupies bin $(14, 23)$; the disordered
A0A0G2L439 occupies $(3, 7)$.
The 42 fragmented proteins cluster in the moderate-to-high arity region.

\begin{figure}[t]
  \centering
  \includegraphics[width=0.9\columnwidth]{figure4_arity_map.png}
  \caption{Arity map of the \textit{H.~sapiens} proteome.
  Cell colour encodes the log-count of proteins with that signature.
  Prototype proteins and fragmented candidates are annotated.}
  \label{fig:arity_map}
\end{figure}

Figure~\ref{fig:q1q4} quantifies the enrichment.
The fragmented set centroid is $(a_{25}=7, a_{75}=13)$, a bin showing a
\textbf{3.08-fold enrichment} relative to the proteome background
(Fisher exact test: $p = 1.76 \times 10^{-7}$).

\begin{figure}[t]
  \centering
  \includegraphics[width=0.9\columnwidth]{figure_q1q4_overlay.png}
  \caption{Q1$\times$Q4 cross-correlation.
  Contours show the arity distribution of all 16{,}787 F1 proteins;
  the yellow ellipse marks the 95\% confidence region of the 42 fragmented
  proteins (centroid at $(7,13)$; 3.08-fold enrichment,
  $p = 1.76 \times 10^{-7}$).}
  \label{fig:q1q4}
\end{figure}

The enrichment reveals a specific structural phenotype: proteins with
moderate $\mathrm{C}_\alpha$ packing density — compact enough to have
a well-defined domain core yet not a single monolithic globule — are
the primary fragmentation candidates.
Highly disordered proteins (low arity, high $H_p$ but PLM $= 1$) and
single-domain globular proteins (high arity, low $H_p$) both fall
outside the fragmented set.
This orthogonal validation from 3D packing geometry supports the
biological interpretation of PLM as a multi-domain indicator.
